% Fichier complété par tools/complete_missing_corrections.py.
% Source : GEO/GEO-03/18_Maxpid/18_Maxpid.tex
% Statut : rédigé (5 question(s)).
\begin{corrigebox}[Corrigé rédigé]

\CorrigeQuestion{1}
Le graphe des liaisons est
$0\xleftrightarrow[\text{pivot}]{D}4\xleftrightarrow[\text{pivot}]{C}3$,
avec $3\xleftrightarrow[\text{glissière}]{\vec x_1}1$,
$1\xleftrightarrow[\text{pivot}]{B}0$ et la liaison hélicoïdale
$2/3$ de pas $p=4\,\mathrm{mm}$. Le rotor 2 est guidé en pivot dans le
stator 1. La chaîne géométrique utile est donc le triangle $BCD$, dont
$BC=\lambda$ et $DC=r$.

\CorrigeQuestion{2}
En notant $L=BD$, $r=DC$ et $\theta_0$ l'angle fixe de
$\overrightarrow{BD}$ avec $\vec x_0$, la fermeture du triangle $BCD$
donne
\[
\lambda^2=L^2+r^2-2Lr\cos(\theta-\theta_0).
\]
Sur la branche correspondant à la configuration dessinée,
\[
\boxed{
\theta(\lambda)=\theta_0+
\arccos\!\left(\frac{L^2+r^2-\lambda^2}{2Lr}\right)
}.
\]

\CorrigeQuestion{3}
La dérivation de la fermeture fournit
\[
2\lambda\dot\lambda
=2Lr\sin(\theta-\theta_0)\dot\theta.
\]
Ainsi, hors position singulière,
\[
\boxed{
\dot\theta=
\frac{\lambda}{Lr\sin(\theta-\theta_0)}\dot\lambda
}.
\]

\CorrigeQuestion{4}
Pour une vis de pas $p$, un tour du rotor produit une translation $p$ :
\[
\dot\lambda=\frac{p}{2\pi}\omega.
\]
Par conséquent,
\[
\boxed{
\dot\theta=
\frac{p\lambda}{2\pi Lr\sin(\theta-\theta_0)}\omega
}.
\]

\CorrigeQuestion{5}
À $n=500\,\mathrm{tr\,min^{-1}}$,
\[
\omega=\frac{2\pi n}{60}=\SI{52.36}{rad.s^{-1}}
\]
et, pour $p=\SI{4}{mm}$,
\[
\dot\lambda=\frac{pn}{60}=\SI{33.3}{mm.s^{-1}}.
\]
Le tracé numérique consiste à balayer les valeurs admissibles de
$\lambda$, puis à calculer successivement
\[
\alpha(\lambda)=
\arccos\!\left(\frac{L^2+r^2-\lambda^2}{2Lr}\right),
\qquad
\theta=\theta_0+\alpha,
\]
et
\[
\dot\theta=
\frac{p\omega}{2\pi}
\frac{\lambda}{Lr\sin\alpha}.
\]
Les deux extrémités du domaine sont écartées, car
$\sin\alpha=0$ dans les configurations singulières.

\end{corrigebox}
